% Assignment source: ne630/homework/html/HW14.html.
% Solution foundation: ne630_problems/lesson_14/hw14_solution.tex.
% Fall 2026 manual formatting pass: removed the standalone-document wrapper
% and reorganized the copied draft into problem/solution blocks consistent
% with the rest of this solution manual.  Problem statements follow the Canvas
% HTML restored to the repository on 2026-10-05.

\noindent{\bf Problem 1}. Consider a very large graphite block with density
$2.26~\mathrm{g\,cm^{-3}}$ and a uniformly distributed source of
$1~\mathrm{MeV}$ neutrons with strength
$s_0~\mathrm{n\,cm^{-3}\,s^{-1}}$.

Using the data supplied in Table~3.2 of FNRP, construct the necessary cross
sections and then compute a three-group approximation of the neutron energy
spectrum.  State all assumptions.

\begin{adaptedsolution}
\noindent\textit{Modeling decisions and available data.}
The problem asks for a coarse spectrum, not a detailed transport calculation,
so the first engineering decision is to use the group structure implied by
FNRP Table~3.2:
\[
\begin{array}{ccl}
g=1 &:& 10^5 < E \le 10^6~\mathrm{eV},\\
g=2 &:& 1 < E \le 10^5~\mathrm{eV},\\
g=3 &:& E \le 1~\mathrm{eV}.
\end{array}
\]
The $1~\mathrm{MeV}$ source is therefore a group-1 source:
\[
S_1=s_0,\qquad S_2=S_3=0.
\]
Let $\phi_g=\int_g\phi(E)\,dE$ be the group-integrated flux.  For isotope
$i$, define microscopic data with the isotope index as a superscript:
\[
\sigma_{a,g}^{(i)}=\sigma_{\gamma,g}^{(i)}+\sigma_{f,g}^{(i)},
\qquad
\Sigma_{x,g}=\sum_i N^{(i)}\sigma_{x,g}^{(i)}10^{-24},
\]
where microscopic cross sections are in barns and atom densities are in
$\mathrm{atoms\,cm^{-3}}$.

The infinite-medium group balance is
\begin{equation}
\Sigma_{r,g}\phi_g
= S_g
+ \sum_{\substack{h=1\\h\ne g}}^3 \Sigma_{s,g\leftarrow h}\phi_h
+ \frac{\chi_g}{k_\infty}s_f''' .
\label{eq:hw14-general-balance}
\end{equation}
For Problem~1 there is no fission source, so the final term is zero.  Removal
from a group counts absorption plus scattering out of that group:
\[
\Sigma_{r,g}=\Sigma_{a,g}
+\sum_{\substack{h=1\\h\ne g}}^3 \Sigma_{s,h\leftarrow g}.
\]

\noindent\textit{Prepare the cross-section data.}
FNRP Table~3.2 gives thermal averages, resonance integrals, and
fission-spectrum fast averages.  The fast values below are therefore
$\chi(E)$-averaged values, not literal $1~\mathrm{MeV}$ point values.  For
the intermediate group, use the resonance-integral weighting
\[
\bar{\sigma}_{x,2}
=\frac{I_x}{\int_1^{10^5}dE/E}
=\frac{I_x}{\ln(10^5)}.
\]
Thus $\Delta u_2=\ln(10^5)=11.512925465$.  The table supplies intermediate
$I_f$ and $I_a$; capture is obtained from
$\sigma_{\gamma,2}=\sigma_{a,2}-\sigma_{f,2}$.  Since Table~3.2 does not
list intermediate scattering averages, use the tabulated thermal scattering
average for group 2.

\begin{center}
\small
\begin{tabular}{llrrrr}
\toprule
Isotope & Group & $\sigma_f$ & $\sigma_\gamma$ & $\sigma_a$ & $\sigma_s$\\
& & \multicolumn{4}{c}{barns}\\
\midrule
\isoA{C}{12} & 1 & 0 & $1.230\times10^{-3}$ & $1.230\times10^{-3}$ & 2.36\\
\isoA{C}{12} & 2 & 0 & $1.329\times10^{-4}$ & $1.329\times10^{-4}$ & 4.81\\
\isoA{C}{12} & 3 & 0 & $3.000\times10^{-3}$ & $3.000\times10^{-3}$ & 4.81\\
\midrule
\isoA{U}{235} & 1 & 1.220 & 0.0700 & 1.290 & 6.33\\
\isoA{U}{235} & 2 & 23.626 & 11.465 & 35.091 & 15.0\\
\isoA{U}{235} & 3 & 505.0 & 86.0 & 591.0 & 15.0\\
\midrule
\isoA{U}{238} & 1 & 0.304 & 0.0570 & 0.361 & 7.42\\
\isoA{U}{238} & 2 & $1.737\times10^{-4}$ & 24.1466 & 24.1468 & 9.37\\
\isoA{U}{238} & 3 & $1.05\times10^{-5}$ & 2.41999 & 2.420 & 9.37\\
\bottomrule
\end{tabular}
\end{center}

The remaining decision is how to map elastic scattering into group-to-group
transfer.  For isotope $i$, define
\[
\alpha^{(i)}=\left(\frac{A^{(i)}-1}{A^{(i)}+1}\right)^2,
\qquad
\xi^{(i)}
=1+\frac{\alpha^{(i)}}{1-\alpha^{(i)}}\ln\alpha^{(i)} .
\]
For elastic scattering with isotropic center-of-mass scattering, the outgoing
energy from incident energy $E$ lies in
$\alpha^{(i)}E\le E'\le E$.  Therefore upscatter terms vanish in this
slowing-down model.  A single elastic collision also cannot carry a group-1
neutron directly to group 3 for these nuclides, so the only off-diagonal
transfer terms are $2\leftarrow1$ and $3\leftarrow2$.  With a $1/E$ spectrum
inside each slowing-down group,
\[
\sigma_{s,g+1\leftarrow g}^{(i)}
\simeq \sigma_{s,g}^{(i)}\frac{\xi^{(i)}}{\Delta u_g},
\qquad
\Delta u_g=\ln\left(\frac{E_{g-1}}{E_g}\right).
\]
Here $\Delta u_1=\ln(10)$ and $\Delta u_2=\ln(10^5)$.

\begin{center}
\small
\begin{tabular}{lrrrr}
\toprule
Isotope & $\alpha^{(i)}$ & $\xi^{(i)}$ &
$\sigma_{s,2\leftarrow1}^{(i)}$ &
$\sigma_{s,3\leftarrow2}^{(i)}$\\
& & & \multicolumn{2}{c}{barns}\\
\midrule
\isoA{C}{12}   & 0.7159763314 & 0.1577689898 & 0.161703 & 0.0659145\\
\isoA{U}{235}  & 0.9831226659 & 0.0084865458 & 0.0233302 & 0.0110570\\
\isoA{U}{238} & 0.9833336251 & 0.0083798718 & 0.0270038 & 0.0068201\\
\bottomrule
\end{tabular}
\end{center}

\noindent\textit{Assemble and solve the graphite balance.}
For pure graphite,
\[
N^{(C)}=\frac{\rho_C}{M_C}N_A
=\frac{2.26}{12.0}(6.02214076\times10^{23})
=1.13416984\times10^{23}~\mathrm{atoms\,cm^{-3}}.
\]
Thus one barn corresponds to
$0.113416984~\mathrm{cm^{-1}}$ for this graphite.  The nonzero macroscopic
cross sections are
\[
\begin{array}{lll}
\Sigma_{a,1}=1.39503\times10^{-4},&
\Sigma_{a,2}=1.50724\times10^{-5},&
\Sigma_{a,3}=3.40251\times10^{-4},\\
\Sigma_{s,2\leftarrow1}=1.83399\times10^{-2},&
\Sigma_{s,3\leftarrow2}=7.47582\times10^{-3}.&
\end{array}
\]
All values are in $\mathrm{cm^{-1}}$.  The fixed-source balances reduce to
\begin{align}
(\Sigma_{a,1}+\Sigma_{s,2\leftarrow1})\phi_1 &= s_0,\\
(\Sigma_{a,2}+\Sigma_{s,3\leftarrow2})\phi_2
  &= \Sigma_{s,2\leftarrow1}\phi_1,\\
\Sigma_{a,3}\phi_3 &= \Sigma_{s,3\leftarrow2}\phi_2.
\label{eq:hw14-fixed-source}
\end{align}
Solving sequentially gives
\[
\phi_1=\frac{s_0}{\Sigma_{a,1}+\Sigma_{s,2\leftarrow1}}
=54.1144\,s_0,
\]
\[
\phi_2=\frac{\Sigma_{s,2\leftarrow1}\phi_1}
{\Sigma_{a,2}+\Sigma_{s,3\leftarrow2}}
=132.488\,s_0,
\]
and
\[
\phi_3=\frac{\Sigma_{s,3\leftarrow2}\phi_2}{\Sigma_{a,3}}
=2910.95\,s_0.
\]
Since $s_0$ has units $\mathrm{n\,cm^{-3}\,s^{-1}}$, these coefficients have
units of cm and the fluxes have units $\mathrm{n\,cm^{-2}\,s^{-1}}$.

\noindent\textit{Interpretation and checks.}
The balance check is
\[
\Sigma_{a,1}\phi_1+\Sigma_{a,2}\phi_2+\Sigma_{a,3}\phi_3=s_0
\]
to rounding.  The result is physically consistent with a nonmultiplying
graphite moderator: most source neutrons scatter down in energy and are
eventually absorbed after reaching the thermal group.
\end{adaptedsolution}

\newpage

\noindent{\bf Problem 2}. Now consider the case in which uranium enriched to
2\% \isoA{U}{235} by atom is uniformly dispersed in the graphite block.
Determine the resulting infinite-medium multiplication factor $k_\infty$ as
a function of the C-to-U atomic ratio $R=N_C/N_U$.  State all assumptions.

\begin{adaptedsolution}
\noindent\textit{Keep the same model, but change the material.}
Use the same group structure, microscopic data, and scattering-transfer model
as in Problem~1.  The material composition is now
\[
N_U=\frac{N_C}{R},\qquad
N^{(235)}=0.02\,\frac{N_C}{R},\qquad
N^{(238)}=0.98\,\frac{N_C}{R}.
\]
Let
\[
c=N_C10^{-24}=0.113416984.
\]
For any microscopic reaction coefficient in uranium, define the enrichment
average
\[
u_{x,g}=0.02\sigma_{x,g}^{(235)}+0.98\sigma_{x,g}^{(238)}.
\]
From the prepared data in Problem~1,
\[
\begin{array}{c|rrr}
g & u_{f,g} & u_{a,g} & u_{\gamma,g}\\
\hline
1 & 0.32232 & 0.37958 & 0.05726\\
2 & 0.47268264 & 24.36565761 & 23.89297497\\
3 & 10.10001029 & 14.1916 & 4.09158971
\end{array}
\]
in barns.  The corresponding uranium-mixture transfer coefficients are
\[
u_{s,2\leftarrow1}=0.02693037,\qquad
u_{s,3\leftarrow2}=0.00690485
\]
in barns.

\noindent\textit{Build the $R$-dependent macroscopic coefficients.}
The absorption and fission coefficients are
\[
\Sigma_{a,g}(R)
=c\left(\sigma_{a,g}^{(C)}+\frac{u_{a,g}}{R}\right),
\qquad
\Sigma_{f,g}(R)=c\frac{u_{f,g}}{R}.
\]
The off-diagonal scattering transfers are
\[
\Sigma_{s,2\leftarrow1}(R)
=c\left(0.161702956+\frac{u_{s,2\leftarrow1}}{R}\right),
\]
\[
\Sigma_{s,3\leftarrow2}(R)
=c\left(0.065914510+\frac{u_{s,3\leftarrow2}}{R}\right).
\]

\noindent\textit{Solve the homogeneous eigenvalue problem.}
For an infinite homogeneous multiplying medium, the source-free group
balances are
\begin{align}
(\Sigma_{a,1}+\Sigma_{s,2\leftarrow1})\phi_1
  &= \frac{s_f'''}{k_\infty},\\
(\Sigma_{a,2}+\Sigma_{s,3\leftarrow2})\phi_2
  &= \Sigma_{s,2\leftarrow1}\phi_1,\\
\Sigma_{a,3}\phi_3 &= \Sigma_{s,3\leftarrow2}\phi_2,
\label{eq:hw14-eig-balances}
\end{align}
where
\[
s_f'''=\sum_{g=1}^3 \nu\Sigma_{f,g}(R)\phi_g .
\]
FNRP Table~3.2 does not give $\bar{\nu}$ values, so the expression below uses
a common named value $\nu$.  If group- or isotope-dependent yields are used,
replace $\nu\Sigma_{f,g}$ by
\[
\sum_i N^{(i)}\bar{\nu}^{(i)}_g\sigma_{f,g}^{(i)}10^{-24}.
\]

Normalize the eigenvector by $\phi_1$:
\[
x(R)=\frac{\phi_2}{\phi_1},\qquad y(R)=\frac{\phi_3}{\phi_1}.
\]
The lower two rows of Eq.~\eqref{eq:hw14-eig-balances} give
\[
x(R)=
\frac{\Sigma_{s,2\leftarrow1}(R)}
{\Sigma_{a,2}(R)+\Sigma_{s,3\leftarrow2}(R)},
\qquad
y(R)=\frac{\Sigma_{s,3\leftarrow2}(R)x(R)}{\Sigma_{a,3}(R)}.
\]
Substitution in the fission production-to-absorption balance gives
\[
\boxed{
k_\infty(R)=
\frac{\nu\Sigma_{f,1}(R)
+\nu\Sigma_{f,2}(R)x(R)+\nu\Sigma_{f,3}(R)y(R)}
{\Sigma_{a,1}(R)+\Sigma_{a,2}(R)x(R)+\Sigma_{a,3}(R)y(R)}
}.
\]

\noindent\textit{Assumptions to attach to this result.}
Fast-group cross sections are fission-spectrum averages rather than pointwise
$1~\mathrm{MeV}$ values; intermediate absorption and fission values are
resonance integrals divided by $\ln(10^5)$; intermediate scattering uses the
thermal scattering averages; elastic transfer is treated as downscatter with
the $1/E$ weighting above; resonance self-shielding is not included; and a
numerical value requires a stated choice of fission neutron yields.
\end{adaptedsolution}

\newpage

\noindent{\bf Problem 3}. Use OpenMC to compute $k_\infty$ for $R=1000$ for
Problem 2 by adapting our in-class example.

\begin{adaptedsolution}
\noindent\textit{Translate the analytic model into an OpenMC input.}
For $R=1000$, the atom densities implied by the graphite density are
\[
N_C=1.13416984\times10^{23}~\mathrm{cm^{-3}},\qquad
N_U=1.13416984\times10^{20}~\mathrm{cm^{-3}},
\]
\[
N^{(235)}=2.26834\times10^{18}~\mathrm{cm^{-3}},\qquad
N^{(238)}=1.11149\times10^{20}~\mathrm{cm^{-3}}.
\]
The OpenMC calculation should therefore use one infinite homogeneous material
containing graphite and 2 atom-percent enriched uranium with $N_C/N_U=1000$.
Use reflective or periodic boundary conditions on a simple cell, run a
$k$-eigenvalue calculation, and report the mean and uncertainty after an
adequate number of inactive and active batches.

\noindent\textit{Analytic check value.}
The three-group model above is not a substitute for OpenMC, but it gives a
useful consistency check.  If the same three-group closure is evaluated with
a common fission neutron yield $\nu=2.43$ for all fissions, then at $R=1000$
\[
x=1.78865,\qquad y=6.85861,
\]
and
\[
\boxed{k_\infty(R=1000;\ \nu=2.43)=1.04793.}
\]
The OpenMC result should be compared against this number as a modeling check,
not treated as a continuous-energy benchmark target.
\end{adaptedsolution}

\clearpage
